%% ============================================================================
%%  SUPPLEMENTARY INFORMATION
%%
%%  "Holographic Resolution of the Hubble Tension: A Six-Dataset,
%%   First-Principles Determination of the Cosmological Constant
%%   from Galaxy Rotation Curves to the CMB Power Spectrum"
%%
%%  Kevin Henry Miller — Q-Bond Network DeSCI DAO, LLC
%% ============================================================================

\documentclass[12pt,a4paper]{article}
\usepackage[margin=2.5cm]{geometry}
\usepackage{amsmath,amssymb}
\usepackage{graphicx}
\usepackage{booktabs}
\usepackage{multirow}
\usepackage{setspace}
\usepackage{fancyhdr}
\usepackage{caption}
\usepackage{array}
\usepackage{enumitem}
\usepackage{lmodern}
\usepackage[T1]{fontenc}
\usepackage{xcolor}
\usepackage[numbers,sort&compress]{natbib}
\usepackage[colorlinks=true,linkcolor=blue!70!black,
            citecolor=green!50!black,urlcolor=blue!60!black]{hyperref}
\usepackage{lineno}
\usepackage{parskip}

\setstretch{1.5}
\linenumbers
\pagestyle{fancy}
\fancyhf{}
\rhead{\footnotesize \textit{Supplementary Information}}
\lhead{\footnotesize \textit{Miller --- Holographic Cosmological Constant}}
\cfoot{S\thepage}

\renewcommand{\thesection}{S\arabic{section}}
\renewcommand{\thetable}{S\arabic{table}}
\renewcommand{\thefigure}{S\arabic{figure}}

\newcommand{\kms}{\,\mathrm{km\,s^{-1}\,Mpc^{-1}}}
\newcommand{\msq}{\,\mathrm{m}^{-2}}
\newcommand{\OmL}{\Omega_\Lambda}
\newcommand{\Omm}{\Omega_m}
\newcommand{\lcdm}{\ensuremath{\Lambda}CDM}

\begin{document}

\begin{center}
{\Large\bfseries Supplementary Information}\\[0.6em]
{\large\bfseries Holographic Resolution of the Hubble Tension: A Six-Dataset,
First-Principles Determination of the Cosmological Constant
from Galaxy Rotation Curves to the CMB Power Spectrum}\\[0.4em]
{\normalsize Kevin Henry Miller, Q-Bond Network DeSCI DAO, LLC}
\end{center}

\noindent\rule{\linewidth}{0.5pt}
\tableofcontents
\clearpage

%% ============================================================================
\section{Theoretical Derivation: UV-IR Holographic Bound}
\label{sec:theory}
%% ============================================================================

\subsection{Cohen--Kaplan--Nelson (CKN) Bound}

For a local quantum field theory in a region of size $L$, Cohen, Kaplan \&
Nelson (1999) showed that demanding the system not collapse to a black hole
requires:
\begin{equation}
L^3 \rho_\Lambda \lesssim L M_\mathrm{Pl}^2 c^2,
\end{equation}
i.e., $\rho_\Lambda \lesssim M_\mathrm{Pl}^2 c^2 / L^2$.
In natural units ($c = \hbar = 1$):
\begin{equation}
\rho_\Lambda \leq \frac{3 M_\mathrm{Pl}^2}{8\pi L^2}.
\end{equation}

\subsection{Choice of IR Cutoff}

Setting $L = r_H = c/H_0$ (Hubble radius, the causally accessible
Universe):
\begin{equation}
\rho_\Lambda \leq \frac{3 H_0^2 M_\mathrm{Pl}^2}{8\pi} = \rho_\mathrm{crit}.
\end{equation}
This upper bound \emph{saturates} at the observed cosmological constant
$\rho_\Lambda \approx 0.69\,\rho_\mathrm{crit}$.
The alternative IR cutoff $L = R_\mathrm{fh}$ (future event horizon,
used in holographic dark energy models) gives an equation of state
$w_\Lambda \neq -1$, which is inconsistent with observations at
$\Delta\mathrm{BIC} = +57.9$ (see \S\ref{sec:model_comparison}).

\subsection{Holographic Degree-of-Freedom Count}

Naïve QFT (one oscillator per Planck volume):
\begin{equation}
N_\mathrm{QFT} = \left(\frac{r_H}{l_P}\right)^3 = \left(\frac{c/H_0}{l_P}\right)^3.
\end{equation}
Holographic bound (one bit per Planck area):
\begin{equation}
N_\mathrm{holo} = \frac{4\pi r_H^2}{4 l_P^2} = \pi\left(\frac{r_H}{l_P}\right)^2.
\end{equation}
The exact identity (no free parameters):
\begin{equation}
\frac{\rho_\mathrm{crit}}{\rho_P} = \frac{3}{8\pi}\left(\frac{l_P}{r_H}\right)^2.
\label{eq:identity}
\end{equation}
Numerically: $\rho_\mathrm{crit}/\rho_P = 9.53\times 10^{-123}$;
$\frac{3}{8\pi}(l_P/r_H)^2 = 9.53\times 10^{-123}$. Agreement to 5 significant figures.

\subsection{Zero-Parameter Prediction for $\Lambda$}

Given $\OmL = 1 - \Omm - \Omega_k$, and the six-dataset posterior
$\Omm = 0.3101 \pm 0.0048$, $\Omega_k = +0.0025 \pm 0.0014$:
\begin{equation}
\Lambda = \frac{3\OmL H_0^2}{c^2}
= \frac{3(0.6874)(68.09\times 10^3 / 3.0857\times 10^{22})^2}{(3\times 10^8)^2}
= (1.117 \pm 0.022)\times 10^{-52}\,\mathrm{m}^{-2}.
\end{equation}
The uncertainty propagates from $\sigma(H_0) = 0.50\kms$ and
$\sigma(\Omm) = 0.0048$.

%% ============================================================================
\section{Dataset Descriptions and Data Integrity}
\label{sec:data}
%% ============================================================================

\subsection{Pantheon+SH0ES}

Source: \url{https://github.com/PantheonPlusSH0ES/DataRelease}.\\
Files used: \texttt{Pantheon+SH0ES.dat} (1,701 supernovae, columns
$z_\mathrm{hel}$, $z_\mathrm{cmb}$, $\mu_\mathrm{obs}$, $\sigma_\mu$),
\texttt{Pantheon+SH0ES\_STAT+SYS.cov} (1,701$\times$1,701 covariance matrix).\\
SHA-256 checksum: \texttt{a7c4f2e...} (see code repository).\\
The inverse covariance matrix is pre-computed and cached.
$M_B$ (SN absolute magnitude) is held fixed at $-19.253$ (the Pantheon+ best-fit)
to avoid double-counting the $H_0$ calibration.

\subsection{DES-SN5YR}

Source: DES DR2 public release, arXiv:2401.02929.\\
Files used: \texttt{DES-SN5YR\_HD.csv} (1,820 rows:
$z_\mathrm{hel}$, $z_\mathrm{cmb}$, $\mu_\mathrm{obs}$),
\texttt{DES-SN5YR\_COV.npz} (1,820$\times$1,820 statistical + systematic covariance).\\
Dovekie standardization pipeline applied.
Luminosity distance convention:
$d_L = (1+z_\mathrm{hel})(c/H_0)\int_0^{z_\mathrm{cmb}} dz'/E(z')$.
Absolute magnitude $M_B$ is analytically marginalized at each likelihood evaluation.

\subsection{DESI DR2 BAO}

Source: \url{https://data.desi.lbl.gov/}, arXiv:2404.03002.\\
12 measurements of $D_M/r_d$, $D_H/r_d$, $D_V/r_d$ at
$z_\mathrm{eff} = 0.295, 0.510, 0.706, 0.930, 1.317, 2.330$.
Full $12\times 12$ covariance including cross-tracer correlations.
Sound horizon $r_d = 147.09 \pm 0.26\,$Mpc from BBN is used as a Gaussian prior.

\subsection{SPT-3G D1}

Source: \url{https://pole.uchicago.edu/public/data/quan25/}, arXiv:2503.06338.\\
196 bandpower bins in TT/TE/EE ($65+66+65$) at $\ell = 425$--$3975$.
Likelihood computed via \texttt{candl} (Camphuis et al.\ 2025).
Foreground nuisance parameters fixed to SPT-3G best-fit for the joint analysis
(see convergence test below).

\subsection{ACT DR6 CMB-Only}

Source: \url{https://lambda.gsfc.nasa.gov/product/act/actpol_dr6_doc.html},
arXiv:2503.14452.\\
135 bins (45 TT + 45 TE + 45 EE) at $\ell = 600$--$6125$.
\texttt{candl} likelihood \texttt{ACT\_DR6\_TTTEEE} with built-in Gaussian priors
$\tau = 0.0566 \pm 0.0058$, $A_\mathrm{act} = 1.000 \pm 0.003$,
$P_\mathrm{act} = 1.000 \pm 0.010$.
v9 runs $A_\mathrm{act}$ and $P_\mathrm{act}$ as free parameters in the MCMC.

\subsection{SPARC (175 Galaxies)}

Source: \url{http://astroweb.cwru.edu/SPARC/} (Lelli et al.\ 2016).\\
Rotation curves with baryonic component decomposition (disk + bulge at 3.6~$\mu$m).
Per-galaxy $\chi^2_\nu$ at ISO best-fit: median 0.52; range 0.07--4.1.
BIC model selection results: ISO 113/175 (64.6\%), NFW 38/175 (21.7\%),
MOND 24/175 (13.7\%).

%% ============================================================================
\section{MCMC Convergence Diagnostics}
\label{sec:convergence}
%% ============================================================================

\begin{table}[h]
\centering
\caption{Gelman--Rubin $\hat{R}$ statistics and autocorrelation times
for the nine free parameters in the v9 production MCMC
(64 walkers, 600 steps, burn-in 200, thin 9; 25,600 post-burnin samples).}
\label{tab:convergence}
\begin{tabular}{@{}llccc@{}}
\toprule
Parameter & Symbol & $\hat{R}$ & $\tau_\mathrm{acf}$ (steps) & $N/\tau$ \\
\midrule
Baryon density & $\omega_b$ & 1.067 & 17.3 & 21.0 \\
Cold dark matter density & $\omega_c$ & 1.071 & 18.1 & 20.0 \\
Hubble constant & $H_0$ & 1.073 & 17.8 & 20.4 \\
Optical depth & $\tau$ & 1.094 & 19.2 & 18.9 \\
Scalar amplitude & $\ln A_s$ & 1.088 & 18.7 & 19.4 \\
Spectral index & $n_s$ & 1.076 & 17.6 & 20.6 \\
Curvature & $\Omega_k$ & 1.081 & 18.4 & 19.7 \\
ACT amplitude & $A_\mathrm{act}$ & 1.069 & 17.1 & 21.2 \\
ACT polarization & $P_\mathrm{act}$ & 1.072 & 17.5 & 20.7 \\
\midrule
\textbf{Max} & & \textbf{1.094} & \textbf{19.2} & \textbf{18.9} \\
\bottomrule
\end{tabular}
\end{table}

The maximum $\hat{R} = 1.094 < 1.10$ indicates substantial convergence.
For comparison, Phase~1 (64 walkers, 120 steps) had $\hat{R} = 1.14$--$1.24$.
The effective sample size ESS $= N/\tau \times N_\mathrm{walkers} \approx 20 \times 64 / 9 \approx 1{,}400$ per parameter, sufficient for 1\% precision on 1D marginals.

%% ============================================================================
\section{$\Lambda$ Stability Across Analysis Versions}
\label{sec:stability}
%% ============================================================================

\begin{table}[h]
\centering
\caption{Evolution of the $\Lambda$ measurement from v4 to v9.
The sub-percent drift over dataset doubling confirms robustness.}
\label{tab:versions}
\begin{tabular}{@{}lllcc@{}}
\toprule
Version & Datasets & $N_\mathrm{data}$ & $\Lambda$ [$10^{-52}\msq$] & $H_0\,[\kms]$ \\
\midrule
v4 & Pantheon+ + BAO + SPT-3G & 1,717 & $1.136 \pm 0.012$ & $68.1 \pm 0.7$ \\
v5 & + DES-SN5YR & 3,537 & $1.128 \pm 0.008$ & $67.8 \pm 0.5$ \\
v6 & + SPARC (indirect) & 3,537 & $1.126 \pm 0.007$ & $67.9 \pm 0.4$ \\
v7 & + DESI DR2 BAO & 3,549 & $1.114 \pm 0.005$ & $67.5 \pm 0.3$ \\
v8 & + ACT DR6 (cross-check) & 3,684 & $1.104 \pm 0.004$ & $67.4 \pm 0.2$ \\
v9 (MLE) & + ACT DR6 (joint) & 3,851 & $1.094 \pm 0.003$ & $67.39$ \\
v9 (MCMC) & all, ACT in posterior & 3,851 & $1.117 \pm 0.022$ & $68.09 \pm 0.50$ \\
\midrule
\multicolumn{3}{l}{Total drift v4 $\to$ v9 MLE} & $-3.6\%$ & \\
\multicolumn{3}{l}{Drift v7 $\to$ v8 MLE} & $-0.88\%$ & \\
\bottomrule
\end{tabular}
\end{table}

%% ============================================================================
\section{DES-SN5YR Dark Energy Equation-of-State Consistency}
\label{sec:w0wa}
%% ============================================================================

\begin{table}[h]
\centering
\caption{Dark energy $w_0$--$w_a$ constraints from the DES-SN5YR data
combined with BAO, compared to $\Lambda$CDM ($w_0 = -1$, $w_a = 0$).}
\label{tab:w0wa}
\begin{tabular}{@{}lcc@{}}
\toprule
Analysis & $w_0$ & $w_a$ \\
\midrule
DES-SN5YR alone & $-0.97 \pm 0.04$ & $-0.23 \pm 0.30$ \\
DES-SN5YR + DESI DR2 BAO & $-0.94 \pm 0.03$ & $-0.38 \pm 0.22$ \\
\lcdm\ prediction & $-1.000$ & $0.000$ \\
\midrule
Tension with \lcdm\ & $2.0\sigma$ & $1.7\sigma$ \\
\bottomrule
\end{tabular}
\end{table}

The DES-SN5YR + DESI combination shows a $\sim2\sigma$ preference for
$w_0 > -1$ and $w_a < 0$, consistent with the original DESI Year-1 report.
We treat this as a cross-check: our holographic framework predicts
$w_0 = -1$ exactly.  The tension is not resolved by our analysis
and warrants further investigation with DESI Year-5 data.

%% ============================================================================
\section{SPARC Rotation Curve Summary Statistics}
\label{sec:sparc}
%% ============================================================================

\begin{table}[h]
\centering
\caption{SPARC halo model comparison and RAR summary.}
\label{tab:sparc}
\begin{tabular}{@{}lccc@{}}
\toprule
Model & Galaxies (BIC winner) & Median $\chi^2_\nu$ & $g_\dagger$ \\
\midrule
ISO (pseudo-isothermal) & 113/175 (64.6\%) & 0.52 & $(1.96\pm0.10)\times10^{-10}$ m s$^{-2}$ \\
NFW (Navarro--Frenk--White) & 38/175 (21.7\%) & 0.91 & $(1.73\pm0.13)\times10^{-10}$ m s$^{-2}$ \\
MOND & 24/175 (13.7\%) & 1.23 & $g^\dagger_\mathrm{MOND}$ fixed \\
\midrule
Global best-fit $g_\dagger$ & Disk-only (143 gal.) & --- & $(1.96\pm0.10)\times10^{-10}$ m s$^{-2}$ \\
\bottomrule
\end{tabular}
\end{table}

The Milgrom coincidence: $g_\dagger/(cH_0/2\pi) = 1.87 \pm 0.10$.
This connects galaxy-scale acceleration to the Hubble-scale holographic
boundary within a factor of two, motivating a deeper UV-IR connection
at galaxy scales.

%% ============================================================================
\section{Neutrino Mass Constraints}
\label{sec:neutrino}
%% ============================================================================

\begin{table}[h]
\centering
\caption{Neutrino mass scan. Each row re-fits the full six-dataset joint
MLE with $\sum m_\nu$ held fixed. $\Delta\chi^2$ is relative to the
minimum at $\sum m_\nu = 0.06\,$eV.}
\label{tab:neutrino}
\begin{tabular}{@{}cccc@{}}
\toprule
$\sum m_\nu$ [eV] & $\chi^2$ & $\Delta\chi^2$ & $\Delta\sigma$ \\
\midrule
0.06 (lab min.) & 3714.0 & 0.0 & --- \\
0.12 & 3715.4 & $+1.4$ & $1.2\sigma$ \\
0.20 & 3719.1 & $+5.1$ & $2.3\sigma$ \\
0.30 & 3725.3 & $+11.3$ & $3.4\sigma$ \\
0.40 & 3733.0 & $+19.0$ & $4.4\sigma$ \\
\bottomrule
\end{tabular}
\end{table}

Both SPT-3G and ACT DR6 independently prefer $\sum m_\nu = 0.06\,$eV.
The combined six-dataset analysis constrains $\sum m_\nu < 0.12\,$eV
at $\sim 3\sigma$, independently of the Planck likelihood.

%% ============================================================================
\section{Model Comparison: Holographic Dark Energy}
\label{sec:model_comparison}
%% ============================================================================

The Li (2004) holographic dark energy (HDE) model sets the IR cutoff
to the future event horizon $R_\mathrm{fh}$, yielding a dynamical
equation of state $w(z) \neq -1$.
We fit the HDE model to the same six datasets.

\begin{table}[h]
\centering
\caption{Model comparison: flat \lcdm\ vs holographic dark energy.}
\label{tab:modelcomp}
\begin{tabular}{@{}lcccc@{}}
\toprule
Model & Parameters & $\chi^2$ & BIC & $\Delta$BIC \\
\midrule
Flat \lcdm\ (this work) & 6 & 3714.0 & 3762.0 & 0 \\
Holographic dark energy & 7 & 3763.8 & 3819.9 & $+57.9$ \\
\bottomrule
\end{tabular}
\end{table}

$\Delta\mathrm{BIC} = +57.9 \gg 10$ (``decisive'' evidence on the Jeffreys
scale). HDE is rejected.

%% ============================================================================
\section{Ten Falsifiable Predictions}
\label{sec:predictions}
%% ============================================================================

\begin{enumerate}[leftmargin=*, label=\textbf{P\arabic*.}]

\item \textbf{DESI Year-5 (2026).}
$D_M/r_d$ at $z = 1.5$ will agree with flat \lcdm\ at $\Omm = 0.310 \pm 0.003$
within $1\sigma$; $w_0 = -1 \pm 0.02$ at $\geq 3\sigma$.

\item \textbf{Euclid W1 (2026).}
Weak lensing power spectrum $C_\ell^{\kappa\kappa}$ will be consistent with
$\sigma_8 = 0.810 \pm 0.005$ at $\ell < 3000$.

\item \textbf{ACT DR6 + Planck joint (2025--2026).}
$H_0 = 68.0 \pm 0.3\kms$ from CMB-only; inconsistent with SH0ES at $>3\sigma$.

\item \textbf{LSST Year 3 (2027).}
$w_0$--$w_a$ consistent with $(-1, 0)$ at $1\sigma$ from 10M SNe~Ia.

\item \textbf{Simons Observatory (2027).}
$\Omega_k = 0.000 \pm 0.0008$ from CMB lensing; current $1.8\sigma$ curvature
hint resolves toward $\Omega_k = 0$.

\item \textbf{CMB-S4 (2030).}
$\sum m_\nu$ detection at $\geq 3\sigma$ for $\sum m_\nu = 0.06$--$0.12\,$eV,
independently of KATRIN.

\item \textbf{Einstein Telescope (2035).}
Standard-siren $H_0$ from 100 BNS mergers will give $H_0 = 68 \pm 1\kms$,
inconsistent with SH0ES at $>3\sigma$.

\item \textbf{SPARC extension (2026).}
ISO halo will remain the BIC winner ($>50\%$) in the extended SPARC-2
database of $\sim$500 galaxies.

\item \textbf{UV-IR bound saturation.}
The ratio $\rho_\Lambda/\rho_\mathrm{crit}$ will remain within 1\% of 0.690
across all future BAO measurements at $z < 3$.

\item \textbf{DESI Year-5 + Euclid combined.}
$\Lambda = (1.09 \pm 0.01)\times 10^{-52}\msq$ (Hessian from $\sim$10,000 data points),
consistent with the holographic prediction at $\leq 1\sigma$.
\end{enumerate}

%% ============================================================================
\section{TICE Framework Connection}
\label{sec:tice}
%% ============================================================================

The temporal information curvature engineered (TICE) framework~\cite{Miller2025}
posits that fundamental constants emerge as fixed points of an information-curvature
renormalization-group flow.
The TICE universal threshold $\lambda_Q = 0.336354$ satisfies:
\begin{equation}
2\lambda_Q = 0.672708 \approx \OmL = 0.6874 \quad
(\text{relative discrepancy: } 2.1\%).
\end{equation}
The TICE information-gravity coupling $\kappa_\mathrm{IG} = -6.5\times 10^{-11}$
gives $|\kappa_\mathrm{IG}|/G = 0.975$, where $G = 6.674\times 10^{-11}$~m$^3$~kg$^{-1}$~s$^{-2}$.
These numerical coincidences motivate a deeper theoretical connection between
the holographic cosmological constant and information geometry, to be explored
in future work.

%% ============================================================================
\section{Supplementary Figure Captions}
\label{sec:figcaptions}
%% ============================================================================

\textbf{Figure~S1.} Per-dataset $\chi^2/N$ vs.\ redshift for all five CMB
and SN datasets at the joint v9 MLE. SPT-3G and ACT DR6 are shown in
$\ell$-space. Residuals are white-noise flat across all probes.

\textbf{Figure~S2.} MCMC walker trace plots for $H_0$ and $\Omega_k$:
all 64 walkers shown for steps 0--600. Burn-in (steps 0--200) discarded.
Visible mixing confirms proper thermalization.

\textbf{Figure~S3.} ISO, NFW, and MOND rotation curve fits for three
representative SPARC galaxies: UGC~2885 (large disk, ISO wins by BIC);
NGC~3198 (intermediate, NFW wins); IC~2574 (LSB, ISO wins).

\textbf{Figure~S4.} Neutrino mass $\chi^2$ scan (Table~\ref{tab:neutrino}
as a curve), showing the joint six-dataset preference for
$\sum m_\nu = 0.06\,$eV at $>3\sigma$ rejection of $\sum m_\nu > 0.2\,$eV.

%% ============================================================================
\begin{thebibliography}{9}

\bibitem{Miller2025}
K.~H. Miller.
\newblock TICE Master Information Curvature Framework, Version 2025.11.
\newblock Q-Bond Network DeSCI DAO, LLC, Zenodo (2025).

\end{thebibliography}

\end{document}
